\chapter{Grundlage}

\section{Grundbegriffe und Definitionen}

Die folgenden Ausführungen führen die zentralen Begriffe ein, die für das Verständnis der im Rahmen dieser Arbeit entwickelten Monitoring-Lösung erforderlich sind. Sie schaffen ein einheitliches Begriffsverständnis als Grundlage für die nachfolgenden Kapitel zur Analyse, Konzeption und Implementierung.

\subsection{Netzwerk-Monitoring, Verfügbarkeit und Latenz}

Unter Netzwerk-Monitoring wird die kontinuierliche Überwachung des Zustands und der Leistungsfähigkeit von Netzwerkkomponenten, Diensten und Schnittstellen verstanden. Dazu gehören beispielsweise Server, Kommunikationsdienste und Webanwendungen, die über ein Netzwerk erreichbar sein müssen. Monitoring-Systeme erfassen hierzu Daten über Verfügbarkeit, Antwortzeiten und Fehlersituationen, um Störungen frühzeitig zu erkennen und geeignete Maßnahmen einleiten zu können \cite{monitoring:prometheus}.

Die Verfügbarkeit eines Dienstes beschreibt, in welchem Anteil einer betrachteten Zeitspanne dieser seinen vorgesehenen Zweck erfüllen kann. Sie wird häufig als Prozentwert angegeben, beispielsweise als \enquote{99{,}9\,\% Verfügbarkeit}. Für die operative Überwachung ist zusätzlich die Erreichbarkeit zu einem konkreten Zeitpunkt relevant. Ein Ziel gilt als erreichbar, wenn eine Verbindung innerhalb eines vorgegebenen Zeitlimits aufgebaut werden kann. Schlägt der Verbindungsaufbau fehl oder wird das Zeitlimit überschritten, wird das Ziel als nicht erreichbar bewertet.

Neben der Erreichbarkeit ist die Latenz eine wichtige Kennzahl zur Beurteilung der Dienstqualität. Sie beschreibt die Zeitspanne zwischen dem Start einer Anfrage und dem Eintreffen einer zugehörigen Antwort beziehungsweise eines Verbindungsstatus. Im Netzwerkumfeld wird sie üblicherweise in Millisekunden angegeben. Eine erhöhte Latenz kann auf Überlastung, Netzwerkprobleme oder eine eingeschränkte Leistungsfähigkeit eines Dienstes hinweisen.

\subsection{Metrikbasiertes Monitoring und Alerting}

Eine Metrik ist eine messbare Kennzahl, die einen bestimmten Aspekt des Zustands oder der Leistungsfähigkeit eines Systems beschreibt. Beispiele sind CPU-Auslastung, Arbeitsspeicherkonsum, Anzahl von Anfragen oder Latenzwerte. Bei einem metriksbasierten Monitoring werden diese Kennzahlen in regelmäßigen Abständen erfasst und als Zeitreihen gespeichert. Dadurch können aktuelle Zustände bewertet sowie Veränderungen über einen längeren Zeitraum analysiert werden \cite{monitoring:prometheus}.

Ein Exporter stellt Metriken in einem standardisierten Format über einen Endpunkt bereit. Ein Monitoring-System kann diese Daten anschließend abholen und speichern. Prometheus verwendet hierfür ein Pull- bzw.\ Scraping-Modell. Dabei ruft der Prometheus-Server die Metrik-Endpunkte der konfigurierten Ziele in festgelegten Intervallen ab \cite{prometheus:configuration}.

Alerting bezeichnet die automatische Erkennung und Meldung kritischer Zustände auf Grundlage von Metriken. Eine Alert-Regel definiert eine Bedingung, bei deren Erfüllung eine Meldung erzeugt wird. Beispiele sind die Nichterreichbarkeit eines Dienstes oder das Überschreiten eines Latenzgrenzwertes. Durch Alerting können Verantwortliche über relevante Störungen informiert werden, ohne ein Dashboard kontinuierlich überwachen zu müssen \cite{prometheus:alerting-overview}.

\subsection{Containerisierung und TCP-basierte Erreichbarkeitsprüfung}

Container sind isolierte Laufzeitumgebungen, in denen Anwendungen zusammen mit ihren Abhängigkeiten ausgeführt werden. Im Vergleich zu virtuellen Maschinen teilen Container den Kernel des Host-Betriebssystems und können daher mit geringerem Ressourcenaufwand gestartet und betrieben werden. Ein Container-Image dient dabei als Vorlage, aus der ein oder mehrere gleichartige Container erzeugt werden können.

Die gemeinsame Verwaltung mehrerer Container wird als Container-Orchestrierung bezeichnet. Sie umfasst unter anderem die Konfiguration von Diensten, Netzwerken, Volumes und Abhängigkeiten. Docker Compose ermöglicht die deklarative Definition und Ausführung von Anwendungen, die aus mehreren Containern bestehen, mithilfe einer YAML-Datei \cite{docker:compose}.

Die TCP-basierte Erreichbarkeitsprüfung dient der Überwachung eines konkreten Dienstes auf einem definierten Host und Port. Das Transmission Control Protocol (TCP) ist verbindungsorientiert und setzt vor einer Datenübertragung einen Verbindungsaufbau zwischen Client und Server voraus. Kann eine TCP-Verbindung erfolgreich aufgebaut werden, ist der Dienst auf dem angegebenen Port grundsätzlich erreichbar. Ein fehlgeschlagener Verbindungsaufbau weist dagegen auf einen nicht verfügbaren Dienst, ein Netzwerkproblem oder eine blockierende Netzwerkkonfiguration hin.

\section{Technologische Grundlagen}

Die entwickelte Monitoring-Lösung basiert auf mehreren miteinander verbundenen Technologien. Spring Boot bildet die Grundlage des Monitoring-Dienstes. Micrometer, Prometheus, Grafana und Alertmanager übernehmen die Erfassung, Speicherung, Visualisierung und Verarbeitung von Monitoring-Daten. Docker und Docker Compose ermöglichen die reproduzierbare Bereitstellung der beteiligten Komponenten.

\subsection{Spring Boot, Actuator und Micrometer}

Spring Boot ist ein Framework zur Entwicklung Java-basierter Anwendungen. Es vereinfacht die Konfiguration typischer Komponenten und unterstützt unter anderem die Bereitstellung von Webschnittstellen und Betriebsfunktionen. Das Actuator-Modul stellt produktionsnahe Funktionen und HTTP-Endpunkte für Zustandsinformationen, Health-Informationen und Metriken bereit \cite{spring-boot-actuator-doc}.

Micrometer ist eine Metrik-Fassade für Java-Anwendungen. Die Bibliothek ermöglicht die Erfassung von Messwerten unabhängig vom konkret eingesetzten Monitoring-Backend. Spring Boot Actuator unterstützt Micrometer durch automatische Konfiguration und kann Metriken für unterschiedliche Monitoring-Systeme bereitstellen \cite{spring-boot-actuator-doc}.

Für Prometheus stellt Spring Boot bei aktivierter Prometheus-Registry den Endpunkt \texttt{/actuator/prometheus} bereit. Dieser Endpunkt liefert Metriken im Prometheus-kompatiblen Format und kann regelmäßig durch einen Prometheus-Server abgefragt werden \cite{spring-boot-actuator-doc}. In der entwickelten Anwendung dient dieser Mechanismus dazu, selbst erzeugte Status- und Latenzmetriken für die weitere Verarbeitung bereitzustellen.

\subsection{Prometheus, Grafana und Alertmanager}

Prometheus ist ein Monitoring-System, das Metriken von konfigurierten Zielen abruft und als Zeitreihen speichert. Die Konfiguration definiert unter anderem Abrufintervalle, Zielsysteme, Metrikpfade und Alert-Regeln. Auf Grundlage der gespeicherten Zeitreihen können Abfragen durchgeführt, Visualisierungen erstellt und kritische Zustände erkannt werden \cite{prometheus:configuration}.

Im entwickelten System fragt Prometheus den Actuator-Endpunkt der Spring-Boot-Anwendung in festen Intervallen ab. Die Anwendung stellt dabei unter anderem die Metriken \texttt{network\_device\_up} und \texttt{network\_device\_latency\_ms} bereit. Diese Werte werden von Prometheus als Zeitreihen gespeichert und können anschließend durch Grafana oder durch Alert-Regeln ausgewertet werden.

Grafana dient der Visualisierung der in Prometheus gespeicherten Daten. Prometheus wird dabei als Datenquelle konfiguriert. Grafana ermöglicht die Darstellung von Zeitreihen, Tabellen und Statusanzeigen, wodurch beispielsweise Verfügbarkeit und Latenz einzelner Netzwerkziele nachvollziehbar werden \cite{grafana:prometheus-datasource}.

Alertmanager verarbeitet Alerts, die von Prometheus aufgrund definierter Regeln erzeugt werden. Zu seinen zentralen Aufgaben zählen die Deduplizierung, Gruppierung und Weiterleitung von Meldungen an konfigurierte Empfänger \cite{prometheus:alertmanager}. In der vorliegenden Lösung können Regeln beispielsweise einen Alert auslösen, wenn ein Ziel über einen definierten Zeitraum nicht erreichbar ist oder ein festgelegter Latenzgrenzwert überschritten wird.

\subsection{Containerisierte Bereitstellung mit Docker und Docker Compose}

Die Monitoring-Lösung wird mit Docker und Docker Compose als mehrteilige Containeranwendung bereitgestellt. Docker Compose ermöglicht es, Services, Netzwerke und Volumes in einer gemeinsamen YAML-Datei zu beschreiben und als zusammengehörigen Anwendungsstack zu starten \cite{docker:compose}.

Die Anwendung besteht aus getrennten Containern für den Spring-Boot-Monitoring-Dienst, Prometheus, Grafana und Alertmanager. Ein gemeinsames Docker-Netzwerk ermöglicht die Kommunikation der Dienste über ihre Servicenamen. Beispielsweise kann Prometheus den Monitoring-Dienst über den Hostnamen \texttt{app} erreichen, während Grafana Prometheus über den Hostnamen \texttt{prometheus} als Datenquelle verwendet.

Für Komponenten mit dauerhaft benötigten Daten, insbesondere Prometheus und Grafana, können Volumes verwendet werden. Damit bleiben Zeitreihendaten und lokale Einstellungen auch nach einem Neustart einzelner Container erhalten. Die containerisierte Bereitstellung reduziert den manuellen Installationsaufwand und erleichtert eine reproduzierbare Ausführung der gesamten Monitoring-Umgebung.

\section{Verwendete Methoden und Konzepte des entwickelten Systems}

Die entwickelte Monitoring-Lösung folgt mehreren methodischen und architektonischen Prinzipien. Im Mittelpunkt stehen eine konfigurationsbasierte Verwaltung der Überwachungsziele, die regelmäßige Überprüfung ihrer Erreichbarkeit sowie die Bereitstellung der Ergebnisse als Metriken. Ergänzend wurden eine klar getrennte Schichtenarchitektur, eine robuste Fehlerbehandlung und automatisierte Tests für zentrale Funktionen umgesetzt.

\subsection{Konfigurationsgetriebenes Gerätemanagement und zeitgesteuertes Polling}

Die zu überwachenden Netzwerkziele werden nicht fest im Quellcode hinterlegt, sondern über die Anwendungskonfiguration definiert. Zu jedem Ziel werden ein Name, ein Host, ein Port sowie ein Aktivierungsstatus angegeben. Zusätzlich werden zentrale Parameter wie das Prüfintervall und das Verbindungs-Timeout konfiguriert.

Die Konfiguration wird beim Start der Anwendung in eine strukturierte Konfigurationsklasse übernommen. Dadurch wird die technische Beschreibung der Überwachungsziele von der eigentlichen Prüfungslogik getrennt. Neue Ziele können durch die Anpassung der Konfigurationsdatei ergänzt werden, ohne dass Änderungen an den Java-Klassen erforderlich sind. Der Aktivierungsstatus ermöglicht außerdem, einzelne Ziele bei Bedarf vorübergehend von der Überwachung auszunehmen.

Die Überwachung erfolgt nach dem Prinzip des zeitgesteuerten Pollings. Dabei überprüft der Monitoring-Dienst alle aktivierten Ziele in regelmäßigen Abständen. Für die Umsetzung wird die Scheduling-Funktion des Spring Framework verwendet. Methoden, die mit der Annotation \texttt{@Scheduled} versehen sind, können in festgelegten Zeitintervallen automatisiert ausgeführt werden \cite{spring:scheduling}.

Das Polling-Verfahren stellt einen Kompromiss zwischen Aktualität und Ressourcenverbrauch dar. Kürzere Intervalle ermöglichen eine schnellere Erkennung von Störungen, verursachen jedoch mehr Verbindungsversuche und Metrikaktualisierungen. Längere Intervalle reduzieren den Aufwand, führen jedoch zu einer größeren Verzögerung bei der Erkennung eines Ausfalls. Durch die Konfigurierbarkeit des Intervalls kann dieser Kompromiss an die jeweilige Einsatzumgebung angepasst werden.

\subsection{Zustandsbasierte Metriken und Kennzahlenmodell}

Die Ergebnisse der Erreichbarkeitsprüfungen werden als zustandsbasierte Metriken modelliert. Für jedes überwachte Ziel werden zwei Kennzahlen bereitgestellt: ein Verfügbarkeitsstatus und eine Latenzmetrik. Der Status wird durch die Metrik \texttt{network\_device\_up} dargestellt. Ein Wert von \texttt{1} bedeutet, dass der Verbindungsaufbau erfolgreich war; ein Wert von \texttt{0} kennzeichnet ein nicht erreichbares Ziel.

Die zweite Metrik \texttt{network\_device\_latency\_ms} enthält die Dauer des letzten Verbindungsversuchs in Millisekunden. Beide Werte werden als Gauges exportiert. Ein Gauge ist für diese Art von Information geeignet, da er einen aktuellen Zustand beschreibt und sich sowohl erhöhen als auch verringern kann. Dies unterscheidet ihn beispielsweise von einem Counter, der ausschließlich monoton ansteigt \cite{prometheus:metric-types}.

Zur Unterscheidung der einzelnen Ziele werden die Metriken mit den Labels \texttt{device} und \texttt{host} versehen. Dadurch können Prometheus-Abfragen, Alert-Regeln und Grafana-Dashboards den Status und die Latenz eines konkreten Ziels auswerten. Das Kennzahlenmodell bildet damit die Grundlage für die Visualisierung und die automatische Erkennung von Störungen.

\subsection{Schichtenarchitektur, Fehlerbehandlung und Robustheit}

Die Anwendung ist in mehrere fachlich getrennte Komponenten gegliedert. Diese Struktur folgt dem Prinzip der Schichtenarchitektur und unterstützt eine klare Zuordnung von Verantwortlichkeiten. Die Konfigurationskomponente übernimmt das Einlesen und Bereitstellen der Einstellungen aus der Datei \texttt{application.yml}. Das Datenmodell umfasst die Klassen \texttt{Device} und \texttt{CheckResult}. Während \texttt{Device} ein zu überwachendes Ziel beschreibt, enthält \texttt{CheckResult} das Ergebnis einer konkreten Prüfung einschließlich Erreichbarkeitsstatus, Latenz und Prüfzeitpunkt.

Die zentrale fachliche Logik befindet sich im \texttt{DeviceCheckService}. Dieser Dienst führt die Prüfungen aus, aktualisiert die Metrikwerte und hält die zuletzt ermittelten Prüfergebnisse bereit. Der \texttt{MonitoringScheduler} stößt die periodische Ausführung an. Die Controller-Schicht stellt REST-Endpunkte bereit, über die Geräte, aktuelle Prüfergebnisse und der Status des Monitoring-Dienstes abgefragt werden können. Zusätzlich ist eine manuelle Auslösung eines Prüfzyklus vorgesehen.

Netzwerkprüfungen sind fehleranfällig, da Verbindungsprobleme, nicht verfügbare Dienste, ungültige Zieladressen oder Zeitüberschreitungen auftreten können. Deshalb werden Ausnahmen beim Aufbau einer TCP-Verbindung behandelt. Schlägt ein Verbindungsaufbau fehl, wird das betreffende Ziel als nicht erreichbar eingestuft; die Prüfung der weiteren Ziele wird dennoch fortgesetzt.

Ein konfigurierbares Timeout begrenzt die Wartezeit für jeden Verbindungsversuch. Dadurch wird verhindert, dass nicht reagierende Ziele den Prüfzyklus über einen längeren Zeitraum blockieren. Auch bei einem fehlgeschlagenen Verbindungsversuch werden Status und Dauer der Prüfung als Ergebnis gespeichert. Die Kombination aus Fehlerbehandlung, Timeout und unabhängiger Verarbeitung einzelner Ziele trägt zur Robustheit der Monitoring-Lösung bei.

\subsection{Automatisierte Tests der Kernlogik}

Zur Absicherung zentraler Funktionen wurden automatisierte Unit-Tests mit JUnit~5 umgesetzt. JUnit stellt einen Rahmen zur Definition und automatisierten Ausführung von Tests für Java-Anwendungen bereit \cite{junit:user-guide}. Der Schwerpunkt liegt auf dem \texttt{DeviceCheckService}, da dieser die Konfiguration der Ziele verarbeitet und die Ergebnisse der Erreichbarkeitsprüfungen erzeugt.

Ein Test prüft, ob die in der Konfiguration definierten Geräte korrekt in die interne Geräteliste übernommen werden. Dadurch wird sichergestellt, dass die konfigurationsgetriebene Verwaltung der Ziele wie vorgesehen funktioniert. Ein weiterer Test verwendet ein nicht erreichbares Ziel und prüft, ob dieses als nicht erreichbar bewertet wird. Zusätzlich wird geprüft, ob für den Verbindungsversuch ein gültiger Latenzwert erzeugt wird.

Die Tests ermöglichen eine wiederholbare Überprüfung der Kernlogik, ohne dass dafür die vollständige Containerumgebung mit Prometheus, Grafana und Alertmanager gestartet werden muss. Sie reduzieren damit das Risiko, dass Änderungen an der Anwendung unbemerkt zu Fehlern in der grundlegenden Monitoring-Funktion führen.

\section{Stand der Technik und Einordnung}

 Blackbox Exporter??? 